\section{Experimental Evaluation}\label{sec:evaluation}
\subsection{Questions, workloads, and evaluation scope}
The case studies address three questions: (1) do software choices change value when hardware changes; (2) what does joint selection gain over fixed or sequential design; and (3) when does the optimization model fail to predict realized behavior? They use a Ventus-derived architecture model with compiled instruction streams, software variants, and parameterized compute, memory, and service resources~\cite{ventus-controlled,ventus-tiling,ventus-model}. The evaluation is organized by these questions rather than by the individual exploration runs.

P128 and P512 denote prefill with 128 and 512 tokens. D128 and D512 denote the mean cycle count over 16 decoding steps starting at the corresponding context length. For each scenario $w$, $r_w=L_w/L_w^{\mathrm{ref}}$ compares a candidate with the same study's fixed baseline. Prefill, decode, and combined objectives take the maximum ratio over their respective scenario sets, as in Equation~\eqref{eq:objective}. A lower ratio is better; a cycle reduction of $1-r$ differs from a speedup of $1/r$.

The controlled fusion study uses Qwen2.5-14B-Instruct with 48 layers, batch one, BF16 operands and storage, and FP32 accumulation. Its 35.52\,mm$^2$ modeled budget uses a fixed ideal resource-replication factor of 32. Whole-network values are calibrated estimates assembled from compiled representative programs, with packing, layout conversion, intermediate movement, bias, and the final language-model head charged. The smaller-network tiling and model-fidelity studies use a 1.11\,mm$^2$ estimated area budget and complete instruction-timing simulations of their selected scenarios. These studies have different evaluators and baselines; their ratios are interpreted within each study.

The evidence evaluates cross-layer mechanisms and bounded MIP-guided selections. The instruction simulator models control, addresses, and resource timing; it does not compute complete inference tensors. Resource and workload-coverage checks accompany the timing results, and numerical checks have a separate local scope. Area is modeled and the results are cycle-based, without a common-frequency physical implementation study. The experiments do not evaluate an autonomous agent producing the entire representation hierarchy. Source snapshots, normalization definitions, and plotted values are provided with the manuscript artifact.

\subsection{Fusion preferences reverse across hardware}\label{sec:fusion-results}
Consider the D128 gate/up projection representative program in the controlled study. On the baseline hardware H0, the fused implementation takes 2.750761 million cycles, compared with 2.918714 million for separate execution. On U3, the separate implementation takes 2.510624 million cycles, while the fused implementation takes 2.762451 million. Fusion is therefore approximately 5.75\% faster on H0; separation is approximately 9.12\% faster on U3. Figure~\ref{fig:controlled}a shows the reversal~\cite{ventus-controlled}.

U3 reallocates the area budget. Relative to H0, it reduces scalar-register slots from 2048 to 128 and local data storage from 128 to 16\,KiB per multiprocessor, while increasing L1 miss-status entries from four to eight and load/store entries from eight to sixteen. Both configurations retain 64 multiprocessors, 1024 vector-register slots, and tensor shape $(4,8,4)$. The measured pair therefore changes more than a scalar throughput parameter: software organization interacts with service concurrency and storage allocation.

This is a direct test of contextual dominance. A rule that permanently removes the separate implementation after evaluating H0 excludes the faster U3 program. The correct knowledge is conditional on the hardware and execution context. A formulation can retain both implementations and link their resource demands and timing coefficients to the selected machine, or a search procedure can reopen the software choice after changing hardware.

\begin{figure}[tbp]
\centering
\begin{tikzpicture}
\begin{axis}[width=6.7cm,height=5.6cm,ybar,bar width=13pt,ymin=0,ymax=3.5,ytick={0,1,2,3},ylabel={Million cycles},symbolic x coords={H0,U3},xtick=data,enlarge x limits=.5,title={(a) Fusion preference reversal},title style={font=\small},tick label style={font=\footnotesize},label style={font=\footnotesize},legend style={font=\scriptsize,at={(.5,-.31)},anchor=north,legend columns=2,draw=none},ymajorgrids,grid style={gray!15}]
\addplot[fill=blueedge!70,draw=blueedge] table[x=hardware,y=fused,col sep=comma] {evidence/fusion.csv};
\addplot[fill=greenedge!65,draw=greenedge] table[x=hardware,y=separate,col sep=comma] {evidence/fusion.csv};
\legend{Fused,Separate}
\end{axis}
\begin{axis}[at={(8.4cm,0)},anchor=south west,width=7.1cm,height=5.6cm,xbar,bar width=5pt,xmin=0,xmax=1.06,xtick={0,.25,.5,.75,1},xlabel={Worst normalized cycles},symbolic y coords={Alternating,Joint,SW then HW,HW only,SW only,Baseline},ytick=data,enlarge y limits=.12,title={(b) Controlled design strategies},title style={font=\small},tick label style={font=\scriptsize},label style={font=\footnotesize},legend style={font=\scriptsize,at={(.5,-.31)},anchor=north,legend columns=2,draw=none},xmajorgrids,grid style={gray!15}]
\addplot[fill=blueedge!70,draw=blueedge] table[x=prefill,y=method,col sep=comma] {evidence/controlled.csv};
\addplot[fill=greenedge!65,draw=greenedge] table[x=decode,y=method,col sep=comma] {evidence/controlled.csv};
\legend{Prefill,Decode}
\end{axis}
\end{tikzpicture}
\caption{Controlled Ventus study~\cite{ventus-controlled}. (a) D128 gate/up representative-program timing changes its preferred software organization across hardware. (b) Whole-network calibrated objectives use the same seven-configuration menu. Joint and iterative alternating selection tie. Lower is better in both panels; the panels use different units.}\label{fig:controlled}
\end{figure}


\subsection{Joint selection and the strength of the baseline}\label{sec:joint-results}
The same study compares baseline, software-only, hardware-only, one-pass software-first, joint, and iterative alternating optimization. The final common menu contains seven hardware configurations; software options include projection fusion, tiling, input packing, attention organization, and row-block dispatch. Each method accesses the same completed menu subject to its own fixed-choice restrictions. Alternating optimization repeatedly computes the best software response to the current hardware and the best hardware response to that software.

For prefill, joint selection reaches $r=0.891559$, a 10.84\% reduction from baseline. For decode, it reaches $r=0.779915$, a 22.01\% reduction. The one-pass software-first objectives are 0.895792 and 0.813232, respectively. Joint selection therefore adds 0.473\% prefill and 4.097\% decode improvement over that sequential procedure. Figure~\ref{fig:controlled}b shows all six strategies. The combined objective selects the same design as prefill~\cite{ventus-controlled}.

The hardware choices are also different across objectives. Prefill selects R072 at an estimated 34.4716\,mm$^2$, while decode selects U3 at 33.3117\,mm$^2$. R072 adds a second writeback port to the U3 organization. On the software side, 32-row dispatch groups improve some prefill projections while retaining the compiled kernel, packed input, and output layout. Paired checks preserve instruction, requested-byte, and matrix-operation counts. Some hardware configurations show little benefit or regress, so the optimizer retains the original dispatch where appropriate.

Iterative alternating optimization reaches the same objectives as joint selection in this menu. The experiment thus establishes the benefit of revisiting coupled decisions, and the loss from a single software-first pass, without attributing an exclusive advantage to the MIP algorithm. This is an informative outcome: when best-response updates suffice, a more elaborate search policy must justify its cost on a larger or more strongly coupled domain.

The selected software and hardware are held fixed for a second-backend evaluation. Prefill and decode reductions remain 10.191\% and 19.941\% relative to that backend's corresponding baselines. Both backends share parts of the instruction-decoding path, so this is a sensitivity check on the selected design. Separate cache-assumption variations produce prefill reductions between 10.66\% and 11.46\%. These ranges describe modeled scenarios rather than statistical confidence intervals.

\subsection{Tiling, registers, and the value of reopening a branch}\label{sec:tiling-results}
A complete-instruction-simulation study provides a complementary example of representation expansion. Its earlier prefill design uses three multiprocessors and tensor shape $(2,16,2)$, with a software mapping based on smaller matrix tiles. A $64\times64$ kernel is then introduced to reduce weight reloads relative to the $32\times64$ tile. The reported traffic analysis identifies repeated loading across query-row blocks, and the larger tile reduces that traffic by roughly one third~\cite{ventus-tiling}.

The new kernel requires 173 vector registers. On the preceding VGPR-512 hardware, the reported comparison worsens from 0.362 to 0.438 in normalized cycles. An alternative hardware organization uses two multiprocessors, tensor shape $(4,8,4)$, VGPR 1024, and a different cache allocation. The larger tile becomes useful in this context. Figure~\ref{fig:trajectory} traces the accepted prefill designs using their complete-simulation objectives.

A second reversal appears within the expanded family. Packed QK, previously unfavorable on the default machine, improves P512 from 0.377 to 0.334 on the new hardware. A subsequent L2 variation produces a design with estimated area 1.0812\,mm$^2$ and ratios 0.3417, 0.3308, 0.2953, and 0.2719 for P128, P512, D128, and D512. The combined objective is 0.3417, corresponding to approximately $2.93\times$ the baseline performance under this timing model. A separate decode-oriented design attains a decode objective near 0.279~\cite{ventus-tiling}.

\begin{figure}[tbp]
\centering
\begin{tikzpicture}
\begin{axis}[width=.89\linewidth,height=5.6cm,ymin=0,ymax=.45,ytick={0,.1,.2,.3,.4},ylabel={Prefill worst normalized cycles},symbolic x coords={Earlier,N3,N4},xtick=data,xticklabels={Earlier accepted design,New tile + hardware,L2 variant},enlarge x limits=.2,ybar,bar width=35pt,nodes near coords,point meta=y,nodes near coords style={font=\footnotesize},every node near coord/.append style={/pgf/number format/fixed,/pgf/number format/precision=4},tick label style={font=\footnotesize},label style={font=\footnotesize},ymajorgrids,grid style={gray!15}]
\addplot[fill=blueedge!65,draw=blueedge] table[x=design,y=prefill,col sep=comma] {evidence/trajectory.csv};
\end{axis}
\end{tikzpicture}
\caption{Accepted-design trajectory from complete instruction simulations. Introducing the larger-tile kernel and a matching register/tensor/cache configuration expands the available combination; an L2 variation then changes the selected hardware. The values are rounded as reported. This trajectory illustrates design evolution, rather than a controlled ablation of individual changes.}\label{fig:trajectory}
\end{figure}


The progression illustrates two operations in the proposed method. Introducing a larger-tile kernel expands the represented program family. Reopening register, tensor, and cache choices creates a coupled neighborhood in which that kernel becomes valuable. It would be misleading to describe this simply as choosing a larger tile or simply as adding registers. The useful object is the new software--hardware combination.

These results are an observed exploration trajectory. They do not isolate the effect of an automated retention policy, since the retained and discarded branches were not compared under a matched policy experiment. They do identify a concrete failure mode that such a policy must handle: a local rejection can become obsolete after a hardware change.

\subsection{A closed MIP gap can coexist with a poor design model}\label{sec:model-results}
The model-fidelity study uses Qwen2.5-0.5B-Instruct, 24 layers, two compiled software families, and a structured search over 27 hardware fields. Nine MIP solves report optimal solutions for their approximate models. Their predicted worst-scenario ratio is approximately 0.3333. Complete instruction simulation instead yields objectives between 0.7913 and 0.8987, a prediction shortfall of roughly 58--63\% relative to the simulated values~\cite{ventus-model}.

The optimization outcome follows the modeling error. Software-only selection retains the baseline software, and hardware-only designs match or outperform the joint selections. For prefill, the joint candidate has objective 0.8987, compared with 0.7913 for hardware-only selection. For the combined objective they tie at 0.8272. An earlier subgraph-selected configuration also outperforms these new candidates on the complete scenarios. The result locates the problem in the connection between model and execution, rather than in an unresolved branch-and-bound gap.

Transfer further exposes the scope of the learned design. Fixing the selected Qwen hardware and applying the corresponding software implementations to GPT-2 small and Pythia-160M produces cycle ratios of 3.07--5.60 relative to their own baselines. In the analyzed cases, the FFN down projection accounts for approximately 62--64\% of the additional cycles, and memory-read volume rises substantially while matched operation instruction counts remain unchanged~\cite{ventus-model}. The observation identifies where to refine the reuse and memory model; it does not isolate one hardware parameter as the cause.

For the proposed control loop, this pattern calls for mechanism and cost revision. Spending more time solving the same approximate MIP cannot correct its ranking. A useful next formulation would expose the mapping-dependent traffic and concurrency responsible for the discrepancy, preserve stronger evaluated designs, and restrict the applicability of any knowledge transferred from the original workload.

\subsection{What the studies establish}
Together, the cases establish three properties of the design problem. Software transformations have hardware-dependent value; a one-pass design process can leave useful combinations unexplored; and model fidelity determines whether mathematical optimization improves the realized design. They also show that sophisticated joint search does not automatically improve on a strong alternating method.

The results support the mechanism-level motivation for progressive co-design. Measuring the full agent policy requires a separate comparison in which candidate access, evaluator, and total exploration budget are controlled. Its outcomes should include design quality, number of costly evaluations, solver effort, and representation changes. The method's benefit would then be attributable to how it constructs and navigates the space, rather than to a larger unreported candidate library.
